\section*{Einsatz generativer KI}
\addcontentsline{toc}{section}{Einsatz generativer KI}

Gemäß den Vorgaben der HTL Leonding zum Einsatz generativer KI sind alle Teile dieser Arbeit, deren Erstfassung von einem Sprachmodell erstellt wurde, im Text mit einer Fußnote gekennzeichnet.
Tabelle~\ref{tab:app:ki} listet diese Teile mit dem verwendeten Werkzeug und dem Link zum Prompt-Verlauf.
Alle generierten Inhalte wurden von den Autoren geprüft, mit Quellen belegt und gegebenenfalls korrigiert.

\reminder{Link zum Chat-Verlauf eintragen (Claude Code: Transkript exportieren und z.\,B. im Repository ablegen) und die letzte Spalte nach der eigenen Überarbeitung anpassen.}

\begin{table}
    \centering
    \small
    \begin{tabularx}{\textwidth}{>{\raggedright\arraybackslash}X l l >{\raggedright\arraybackslash}X}
        \toprule
        \textbf{Teil der Arbeit} & \textbf{Werkzeug} & \textbf{Datum} & \textbf{Prompt-Verlauf} \\
        \midrule
        Abstract, Zusammenfassung (Kurzfassung) & Claude Opus 5.5 & 13.10.2026 & TODO Link \\
        Kapitel~\ref{chapter:introduction} Einleitung & Claude Opus 5.5 & 13.10.2026 & TODO Link \\
        Kapitel~\ref{chapter:related} Umfeldanalyse & Claude Opus 5.5 & 13.10.2026 & TODO Link \\
        Kapitel~\ref{chapter:requirements} Anforderungen & Claude Opus 5.5 & 13.10.2026 & TODO Link \\
        Kapitel~\ref{chapter:architecture} Systemüberblick & Claude Opus 5.5 & 13.10.2026 & TODO Link \\
        \bottomrule
    \end{tabularx}
    \caption{Mit generativer KI erstellte Erstfassungen}
    \label{tab:app:ki}
\end{table}
